\src{171}
\result{PROPOSITION}{3} Let $L_{*}=(L_i)$ be a left resolution of the
$\Osh$-Module $A$ by $\Osh$-Modules, each locally isomorphic to some $\Osh^{n}$.
Then $\Extsh_{\Osh}^{*}(A,B)$ identifies with
$H^{*}\bigl(\Homsh_{\Osh}(L_{*},B)\bigr)$, and
$\operatorname{Ext}_{\Osh}^{*}(X;A,B)$ identifies with the hypercohomology of
$X$ with respect to the complex $\Homsh_{\Osh}(L_{*},B)$.

The proof is standard: consider the bicomplex $\Homsh_{\Osh}(L_{*},C(B))$,
where $C(B)$ is an injective resolution of $B$, and the natural homomorphisms
from $\Homsh_{\Osh}(L_{*},B)$ and $\Homsh_{\Osh}(A,C(B))$ into this bicomplex.

Finally, let us point out that the two $\operatorname{Ext}$ functors introduced
in (1.2) are not only cohomological functors in $B$, but cohomological
bifunctors, covariant in $B$ and contravariant in $A$.

\fsection{2}{Composition law in Ext}

The results of this section are due independently to CARTIER and YONEDA; see an
exposé by CARTIER [1] for details. Let $\Ccat$ be an abelian category, and let
$K$ and $L$ be two graded objects in $\Ccat$. We denote by
$\operatorname{Hom}(K,L)$ the graded abelian group whose component of degree
$n$ consists of the homogeneous homomorphisms of degree $n$ from $K$ to $L$
(that is, systems $(u_i)$ of homomorphisms
$K^{i}\longrightarrow L^{i+n}$).

If $K$ and $L$ are complexes (with differential operators of degree $+1$, to
fix ideas), we equip $\operatorname{Hom}(K,L)$ with the differential operator
given by
\begin{equation}\tag{2.1}
  \delta(u)=du+(-1)^{n+1}ud,\qquad\text{where }n=\deg(u),
\end{equation}
making it a complex with differential of degree $+1$. The cycles of degree $n$
are the operators of degree $n$ that anticommute with $d$ (as homogeneous
operators).\footnote{The printed text reads “with $u$”; equation (2.1) requires
“with $d$.”} We may therefore consider
$H^{*}\bigl(\operatorname{Hom}(K,L)\bigr)$; it is an invariant of the homotopy
types of $K$ and $L$, and may be denoted by $H^{*}(K,L)$.

Given a third complex $M$, composition of homomorphisms defines a pairing
\[
  \operatorname{Hom}(K,L)\times\operatorname{Hom}(L,M)
  \longrightarrow\operatorname{Hom}(K,M)
\]
compatible with the differentials, hence, after passing to cohomology, pairings
\begin{equation}\tag{2.2}
  H^{*}(K,L)\times H^{*}(L,M)\longrightarrow H^{*}(K,M),
\end{equation}
denoted $(u,v)\longmapsto vu$. These pairings satisfy the evident associativity
property. In particular, $H^{*}(K,K)$ is an associative graded ring with
identity, and $H^{*}(K,L)$ (respectively $H^{*}(L,K)$) is a graded right
(respectively left) module over this ring, and so on. In degree $0$, (2.2)
reduces to composition of morphisms of complexes. Finally, an exact sequence of
complexes
\[
  0\longrightarrow K'\longrightarrow K\longrightarrow K''\longrightarrow 0,
\]
